% !TeX program = lualatex
% !TeX root = predavanje09.tex
% ============================================================
% Zajednički preamble za sva predavanja RWA kolegija
% Uključiti s: % ============================================================
% Zajednički preamble za sva predavanja RWA kolegija
% Uključiti s: % ============================================================
% Zajednički preamble za sva predavanja RWA kolegija
% Uključiti s: \input{../common/preamble.tex}
% ============================================================

\documentclass[aspectratio=149,11pt]{beamer}
\usetheme{CambridgeUS}

\usepackage{fontspec}              % radi samo sa XeLaTeX/LuaLaTeX
\usepackage[croatian]{babel}
\usepackage{csquotes}              % \enquote{} → hrvatski navodnici „..."
\usepackage{booktabs}
\usepackage{array}
\usepackage{listings}
\usepackage{xcolor}
\usepackage{colortbl}
\usepackage{hyperref}
\usepackage{tikz}
\usetikzlibrary{positioning,arrows.meta}

% ---------- Tipografija: IBM Plex ----------
\setsansfont{IBMPlexSans}[
  Extension      = .otf,
  UprightFont    = *-Regular,
  ItalicFont     = *-Italic,
  BoldFont       = *-SemiBold,
  BoldItalicFont = *-SemiBoldItalic
]
\setmainfont{IBMPlexSans}[
  Extension      = .otf,
  UprightFont    = *-Regular,
  ItalicFont     = *-Italic,
  BoldFont       = *-SemiBold,
  BoldItalicFont = *-SemiBoldItalic
]
\setmonofont{IBMPlexMono}[
  Extension      = .otf,
  UprightFont    = *-Regular,
  ItalicFont     = *-Italic,
  BoldFont       = *-SemiBold,
  BoldItalicFont = *-SemiBoldItalic,
  Scale          = 0.92
]

% ---------- Paleta ----------
\definecolor{ink}{HTML}{0C4685}      % primarna tamnoplava
\definecolor{ink2}{HTML}{6E7BAA}     % sekundarna plavo-siva
\definecolor{tint}{HTML}{EAEEF6}     % svjetli plavi tint
\definecolor{textclr}{HTML}{1A222E}  % glavni tekst
\definecolor{mutedclr}{HTML}{6B7484} % podnožje, sekundarni tekst
\definecolor{bgclr}{HTML}{F6F7F9}    % pozadina slajda
\definecolor{lineclr}{HTML}{D9DEE7}  % tanke linije tablica
\definecolor{divrule}{HTML}{9DB4D6}  % akcent na tamnim sekcijskim slajdovima
\definecolor{divtext}{HTML}{E6ECF6}  % tekst na tamnim sekcijskim slajdovima

\setbeamercolor{normal text}{fg=textclr,bg=bgclr}
\setbeamercolor{structure}{fg=ink}
\setbeamercolor{alerted text}{fg=ink}

\setbeamercolor{title}{fg=ink,bg=tint}
\setbeamercolor{subtitle}{fg=mutedclr}
\setbeamercolor{author}{fg=textclr}
\setbeamercolor{date}{fg=mutedclr}
\setbeamercolor{frametitle}{fg=ink,bg=tint}

\setbeamercolor{section in head/foot}{fg=white,bg=ink}
\setbeamercolor{subsection in head/foot}{fg=ink,bg=tint}
\setbeamercolor{author in head/foot}{fg=white,bg=ink}
\setbeamercolor{title in head/foot}{fg=textclr,bg=tint}
\setbeamercolor{date in head/foot}{fg=mutedclr,bg=tint}

\setbeamercolor{block title}{fg=white,bg=ink}
\setbeamercolor{block body}{fg=textclr,bg=tint}
\setbeamercolor{block title example}{fg=white,bg=ink2}
\setbeamercolor{block body example}{fg=textclr,bg=tint}
\setbeamercolor{itemize item}{fg=ink2}
\setbeamercolor{itemize subitem}{fg=ink2}
\setbeamercolor{enumerate item}{fg=ink2}

\setbeamertemplate{itemize items}[square]
\setbeamertemplate{navigation symbols}{}
\arrayrulecolor{lineclr}

% ---------- Kolegij, autori i logo ----------
\newcommand{\kolegij}{Programiranje za web}
\author{doc.\,dr.\,sc. Tomislav Petković \and mag.\,ing. Josip Torić}
\newcommand{\znak}[2]{\raisebox{-0.5\height}{\includegraphics[height=#1]{#2}}}
\titlegraphic{\znak{2.1cm}{../common/SIT_logo_tekst1.pdf}\hspace{1.8em}\znak{1.8cm}{../common/UNIZD_logo.pdf}}

% ---------- Zaglavlje ----------
% Jedna puna traka: lijevo sekcija, desno podsekcija.
\setbeamertemplate{headline}{%
  \leavevmode%
  \hbox{%
    \begin{beamercolorbox}[wd=\paperwidth,ht=2.25ex,dp=1ex,leftskip=.3cm,rightskip=.3cm]{section in head/foot}%
      \usebeamerfont{section in head/foot}\insertsectionhead\hfill\insertsubsectionhead%
    \end{beamercolorbox}%
  }%
}

% ---------- Podnožje ----------
% Lijevo: naziv trenutne prezentacije (kratki oblik naslova iz \title[...]{...})
% Sredina: naziv kolegija.  Desno: datum i redni broj slajda.
\setbeamertemplate{footline}{%
  \leavevmode%
  \hbox{%
    \begin{beamercolorbox}[wd=.34\paperwidth,ht=2.25ex,dp=1ex,center]{author in head/foot}%
      \usebeamerfont{author in head/foot}\insertshorttitle%
    \end{beamercolorbox}%
    \begin{beamercolorbox}[wd=.33\paperwidth,ht=2.25ex,dp=1ex,center]{title in head/foot}%
      \usebeamerfont{title in head/foot}\kolegij%
    \end{beamercolorbox}%
    \begin{beamercolorbox}[wd=.33\paperwidth,ht=2.25ex,dp=1ex,center]{date in head/foot}%
      \usebeamerfont{date in head/foot}\insertshortdate{}\hspace*{2em}%
      \insertframenumber{}\,/\,\inserttotalframenumber\hspace*{1ex}%
    \end{beamercolorbox}%
  }%
  \vskip0pt%
}

% ---------- Sekcijski slajd ----------
% \sekcija{Naslov sekcije}{jedna linija koja nabraja teme sekcije}
\newcommand{\sekcija}[2]{%
  \section{#1}%
  {%
    \setbeamercolor{background canvas}{bg=ink}%
    \setbeamertemplate{headline}{}%
    \begin{frame}
      \vfill
      \hspace*{0.03\textwidth}{\color{divrule}\rule{2.6cm}{1.1pt}}\par
      \vskip 0.7em
      \hspace*{0.03\textwidth}{\Large\bfseries\color{white}#1}\par
      \vskip 0.6em
      \hspace*{0.03\textwidth}\begin{minipage}{0.9\textwidth}\color{divtext}\small #2\end{minipage}\par
      \vfill
    \end{frame}%
  }%
}

% ---------- Listings ----------
\lstdefinelanguage{bash}{
  morekeywords={sudo,apt,brew,pip,pipx,uv,winget,python,git,flask,source,deactivate,code,gunicorn,docker,compose,heroku,render,railway,node,npm,pnpm,yarn,uvicorn,alembic,pytest,ruff,black,precommit},
  sensitive=false,
  morecomment=[l]{\#}
}
\lstdefinelanguage{python}{
  morekeywords={from,import,def,class,return,if,elif,else,for,in,while,try,except,raise,with,as,not,and,or,is,lambda,yield,True,False,None,await,async},
  sensitive=true,
  morecomment=[l]{\#},
  morestring=[b]',
  morestring=[b]"
}
\lstdefinelanguage{sql}{
  morekeywords={SELECT,FROM,WHERE,INSERT,INTO,VALUES,UPDATE,SET,DELETE,CREATE,TABLE,PRIMARY,KEY,FOREIGN,REFERENCES,UNIQUE,NOT,NULL,INDEX,ON,ALTER,DROP,CONSTRAINT},
  sensitive=false,
  morecomment=[l]{--}
}
\lstdefinelanguage{json}{
  morecomment=[l]{//},
  morestring=[b]",
  showstringspaces=false
}
\lstdefinelanguage{typescript}{
  morekeywords={import,export,from,const,let,var,function,return,if,else,for,of,in,while,switch,case,break,async,await,try,catch,finally,throw,new,class,this,extends,implements,interface,type,enum,readonly,public,private,protected,static,get,set,true,false,null,undefined,void,string,number,boolean,any,unknown,never,typeof,keyof,as},
  sensitive=true,
  morecomment=[l]{//},
  morecomment=[s]{/*}{*/},
  morestring=[b]',
  morestring=[b]",
  morestring=[b]`
}
\lstalias{javascript}{typescript}
\lstset{
  backgroundcolor=\color{tint},
  basicstyle=\ttfamily\footnotesize,
  keywordstyle=\color{ink}\bfseries,
  commentstyle=\color{mutedclr},
  frame=single,
  rulecolor=\color{lineclr},
  frameround=tttt,
  breaklines=true,
  tabsize=2,
  showstringspaces=false,
  upquote=true,
  columns=fullflexible,
  escapeinside={<@}{@>}
}
% ---------- Dijakritici u isječcima koda ----------
% listings čita bajt po bajt, pa višebajtni znak premjesti: "početni" ispadne
% "čpoetni". literate to ne popravlja. Rješenje je izlaz u LaTeX: sve između
% <@ i @> ne prolazi kroz listings, nego se slaže kao običan tekst, a unutar
% okvira je font ionako ttfamily, pa se ništa ne razlikuje.
%
%   async def get(...):
%       club = await repo.get_by_id(db, club_id)   <@\kom{dohvat kluba}@>
%       raise AppError("not_found", <@"Klub nije pronađen"@>, 404)
%
% Graničnici <@ i @> odabrani su jer se ne pojavljuju ni u jednom jeziku
% koji se na predavanjima citira; (* bi se sudaralo s def f(*args).
\newcommand{\kom}[1]{{\color{mutedclr}\# #1}}
% Isto za TypeScript, Vue i CSS, kojima komentar počinje s //:
%   const lifters = useLiftersStore()   <@\komts{dijeljeni popis}@>
\newcommand{\komts}[1]{{\color{mutedclr}// #1}}

% ---------- Činjenica naspram konvencije ----------
% Sve što nije propisano standardom ili jezikom, nego je dogovor struke,
% ide u ovaj okvir, da student zna koje je vrste tvrdnja pred njim.
% \begin{konvencija} ... \end{konvencija}
% \begin{konvencija}[Konvencija: nazivi putanja] ... \end{konvencija}
\newenvironment{konvencija}[1][Konvencija]
  {\begin{exampleblock}{#1}}
  {\end{exampleblock}}

% ---------- Helper macros ----------
% Stupac tablice zadane širine, poravnat ulijevo (bez rastezanja i dijeljenja riječi).
\newcolumntype{L}[1]{>{\raggedright\arraybackslash}p{#1}}
\newcommand{\codefile}[1]{\texttt{#1}}
\newcommand{\kbd}[1]{\texttt{#1}}

% ---------- Code frame environment ----------
% Automatski dodaje fragile i shrink=5, koristi se umjesto \begin{frame}
% Primjer: \begin{codeframe}{Naslov} ... \end{codeframe}
% Primjer s opcijama: \begin{codeframe}[shrink=15]{Naslov} ... \end{codeframe}
\newenvironment{codeframe}[2][]
  {\begin{frame}[fragile,shrink=5,#1]{#2}}
  {\end{frame}}

% ============================================================

\documentclass[aspectratio=149,11pt]{beamer}
\usetheme{CambridgeUS}

\usepackage{fontspec}              % radi samo sa XeLaTeX/LuaLaTeX
\usepackage[croatian]{babel}
\usepackage{csquotes}              % \enquote{} → hrvatski navodnici „..."
\usepackage{booktabs}
\usepackage{array}
\usepackage{listings}
\usepackage{xcolor}
\usepackage{colortbl}
\usepackage{hyperref}
\usepackage{tikz}
\usetikzlibrary{positioning,arrows.meta}

% ---------- Tipografija: IBM Plex ----------
\setsansfont{IBMPlexSans}[
  Extension      = .otf,
  UprightFont    = *-Regular,
  ItalicFont     = *-Italic,
  BoldFont       = *-SemiBold,
  BoldItalicFont = *-SemiBoldItalic
]
\setmainfont{IBMPlexSans}[
  Extension      = .otf,
  UprightFont    = *-Regular,
  ItalicFont     = *-Italic,
  BoldFont       = *-SemiBold,
  BoldItalicFont = *-SemiBoldItalic
]
\setmonofont{IBMPlexMono}[
  Extension      = .otf,
  UprightFont    = *-Regular,
  ItalicFont     = *-Italic,
  BoldFont       = *-SemiBold,
  BoldItalicFont = *-SemiBoldItalic,
  Scale          = 0.92
]

% ---------- Paleta ----------
\definecolor{ink}{HTML}{0C4685}      % primarna tamnoplava
\definecolor{ink2}{HTML}{6E7BAA}     % sekundarna plavo-siva
\definecolor{tint}{HTML}{EAEEF6}     % svjetli plavi tint
\definecolor{textclr}{HTML}{1A222E}  % glavni tekst
\definecolor{mutedclr}{HTML}{6B7484} % podnožje, sekundarni tekst
\definecolor{bgclr}{HTML}{F6F7F9}    % pozadina slajda
\definecolor{lineclr}{HTML}{D9DEE7}  % tanke linije tablica
\definecolor{divrule}{HTML}{9DB4D6}  % akcent na tamnim sekcijskim slajdovima
\definecolor{divtext}{HTML}{E6ECF6}  % tekst na tamnim sekcijskim slajdovima

\setbeamercolor{normal text}{fg=textclr,bg=bgclr}
\setbeamercolor{structure}{fg=ink}
\setbeamercolor{alerted text}{fg=ink}

\setbeamercolor{title}{fg=ink,bg=tint}
\setbeamercolor{subtitle}{fg=mutedclr}
\setbeamercolor{author}{fg=textclr}
\setbeamercolor{date}{fg=mutedclr}
\setbeamercolor{frametitle}{fg=ink,bg=tint}

\setbeamercolor{section in head/foot}{fg=white,bg=ink}
\setbeamercolor{subsection in head/foot}{fg=ink,bg=tint}
\setbeamercolor{author in head/foot}{fg=white,bg=ink}
\setbeamercolor{title in head/foot}{fg=textclr,bg=tint}
\setbeamercolor{date in head/foot}{fg=mutedclr,bg=tint}

\setbeamercolor{block title}{fg=white,bg=ink}
\setbeamercolor{block body}{fg=textclr,bg=tint}
\setbeamercolor{block title example}{fg=white,bg=ink2}
\setbeamercolor{block body example}{fg=textclr,bg=tint}
\setbeamercolor{itemize item}{fg=ink2}
\setbeamercolor{itemize subitem}{fg=ink2}
\setbeamercolor{enumerate item}{fg=ink2}

\setbeamertemplate{itemize items}[square]
\setbeamertemplate{navigation symbols}{}
\arrayrulecolor{lineclr}

% ---------- Kolegij, autori i logo ----------
\newcommand{\kolegij}{Programiranje za web}
\author{doc.\,dr.\,sc. Tomislav Petković \and mag.\,ing. Josip Torić}
\newcommand{\znak}[2]{\raisebox{-0.5\height}{\includegraphics[height=#1]{#2}}}
\titlegraphic{\znak{2.1cm}{../common/SIT_logo_tekst1.pdf}\hspace{1.8em}\znak{1.8cm}{../common/UNIZD_logo.pdf}}

% ---------- Zaglavlje ----------
% Jedna puna traka: lijevo sekcija, desno podsekcija.
\setbeamertemplate{headline}{%
  \leavevmode%
  \hbox{%
    \begin{beamercolorbox}[wd=\paperwidth,ht=2.25ex,dp=1ex,leftskip=.3cm,rightskip=.3cm]{section in head/foot}%
      \usebeamerfont{section in head/foot}\insertsectionhead\hfill\insertsubsectionhead%
    \end{beamercolorbox}%
  }%
}

% ---------- Podnožje ----------
% Lijevo: naziv trenutne prezentacije (kratki oblik naslova iz \title[...]{...})
% Sredina: naziv kolegija.  Desno: datum i redni broj slajda.
\setbeamertemplate{footline}{%
  \leavevmode%
  \hbox{%
    \begin{beamercolorbox}[wd=.34\paperwidth,ht=2.25ex,dp=1ex,center]{author in head/foot}%
      \usebeamerfont{author in head/foot}\insertshorttitle%
    \end{beamercolorbox}%
    \begin{beamercolorbox}[wd=.33\paperwidth,ht=2.25ex,dp=1ex,center]{title in head/foot}%
      \usebeamerfont{title in head/foot}\kolegij%
    \end{beamercolorbox}%
    \begin{beamercolorbox}[wd=.33\paperwidth,ht=2.25ex,dp=1ex,center]{date in head/foot}%
      \usebeamerfont{date in head/foot}\insertshortdate{}\hspace*{2em}%
      \insertframenumber{}\,/\,\inserttotalframenumber\hspace*{1ex}%
    \end{beamercolorbox}%
  }%
  \vskip0pt%
}

% ---------- Sekcijski slajd ----------
% \sekcija{Naslov sekcije}{jedna linija koja nabraja teme sekcije}
\newcommand{\sekcija}[2]{%
  \section{#1}%
  {%
    \setbeamercolor{background canvas}{bg=ink}%
    \setbeamertemplate{headline}{}%
    \begin{frame}
      \vfill
      \hspace*{0.03\textwidth}{\color{divrule}\rule{2.6cm}{1.1pt}}\par
      \vskip 0.7em
      \hspace*{0.03\textwidth}{\Large\bfseries\color{white}#1}\par
      \vskip 0.6em
      \hspace*{0.03\textwidth}\begin{minipage}{0.9\textwidth}\color{divtext}\small #2\end{minipage}\par
      \vfill
    \end{frame}%
  }%
}

% ---------- Listings ----------
\lstdefinelanguage{bash}{
  morekeywords={sudo,apt,brew,pip,pipx,uv,winget,python,git,flask,source,deactivate,code,gunicorn,docker,compose,heroku,render,railway,node,npm,pnpm,yarn,uvicorn,alembic,pytest,ruff,black,precommit},
  sensitive=false,
  morecomment=[l]{\#}
}
\lstdefinelanguage{python}{
  morekeywords={from,import,def,class,return,if,elif,else,for,in,while,try,except,raise,with,as,not,and,or,is,lambda,yield,True,False,None,await,async},
  sensitive=true,
  morecomment=[l]{\#},
  morestring=[b]',
  morestring=[b]"
}
\lstdefinelanguage{sql}{
  morekeywords={SELECT,FROM,WHERE,INSERT,INTO,VALUES,UPDATE,SET,DELETE,CREATE,TABLE,PRIMARY,KEY,FOREIGN,REFERENCES,UNIQUE,NOT,NULL,INDEX,ON,ALTER,DROP,CONSTRAINT},
  sensitive=false,
  morecomment=[l]{--}
}
\lstdefinelanguage{json}{
  morecomment=[l]{//},
  morestring=[b]",
  showstringspaces=false
}
\lstdefinelanguage{typescript}{
  morekeywords={import,export,from,const,let,var,function,return,if,else,for,of,in,while,switch,case,break,async,await,try,catch,finally,throw,new,class,this,extends,implements,interface,type,enum,readonly,public,private,protected,static,get,set,true,false,null,undefined,void,string,number,boolean,any,unknown,never,typeof,keyof,as},
  sensitive=true,
  morecomment=[l]{//},
  morecomment=[s]{/*}{*/},
  morestring=[b]',
  morestring=[b]",
  morestring=[b]`
}
\lstalias{javascript}{typescript}
\lstset{
  backgroundcolor=\color{tint},
  basicstyle=\ttfamily\footnotesize,
  keywordstyle=\color{ink}\bfseries,
  commentstyle=\color{mutedclr},
  frame=single,
  rulecolor=\color{lineclr},
  frameround=tttt,
  breaklines=true,
  tabsize=2,
  showstringspaces=false,
  upquote=true,
  columns=fullflexible,
  escapeinside={<@}{@>}
}
% ---------- Dijakritici u isječcima koda ----------
% listings čita bajt po bajt, pa višebajtni znak premjesti: "početni" ispadne
% "čpoetni". literate to ne popravlja. Rješenje je izlaz u LaTeX: sve između
% <@ i @> ne prolazi kroz listings, nego se slaže kao običan tekst, a unutar
% okvira je font ionako ttfamily, pa se ništa ne razlikuje.
%
%   async def get(...):
%       club = await repo.get_by_id(db, club_id)   <@\kom{dohvat kluba}@>
%       raise AppError("not_found", <@"Klub nije pronađen"@>, 404)
%
% Graničnici <@ i @> odabrani su jer se ne pojavljuju ni u jednom jeziku
% koji se na predavanjima citira; (* bi se sudaralo s def f(*args).
\newcommand{\kom}[1]{{\color{mutedclr}\# #1}}
% Isto za TypeScript, Vue i CSS, kojima komentar počinje s //:
%   const lifters = useLiftersStore()   <@\komts{dijeljeni popis}@>
\newcommand{\komts}[1]{{\color{mutedclr}// #1}}

% ---------- Činjenica naspram konvencije ----------
% Sve što nije propisano standardom ili jezikom, nego je dogovor struke,
% ide u ovaj okvir, da student zna koje je vrste tvrdnja pred njim.
% \begin{konvencija} ... \end{konvencija}
% \begin{konvencija}[Konvencija: nazivi putanja] ... \end{konvencija}
\newenvironment{konvencija}[1][Konvencija]
  {\begin{exampleblock}{#1}}
  {\end{exampleblock}}

% ---------- Helper macros ----------
% Stupac tablice zadane širine, poravnat ulijevo (bez rastezanja i dijeljenja riječi).
\newcolumntype{L}[1]{>{\raggedright\arraybackslash}p{#1}}
\newcommand{\codefile}[1]{\texttt{#1}}
\newcommand{\kbd}[1]{\texttt{#1}}

% ---------- Code frame environment ----------
% Automatski dodaje fragile i shrink=5, koristi se umjesto \begin{frame}
% Primjer: \begin{codeframe}{Naslov} ... \end{codeframe}
% Primjer s opcijama: \begin{codeframe}[shrink=15]{Naslov} ... \end{codeframe}
\newenvironment{codeframe}[2][]
  {\begin{frame}[fragile,shrink=5,#1]{#2}}
  {\end{frame}}

% ============================================================

\documentclass[aspectratio=149,11pt]{beamer}
\usetheme{CambridgeUS}

\usepackage{fontspec}              % radi samo sa XeLaTeX/LuaLaTeX
\usepackage[croatian]{babel}
\usepackage{csquotes}              % \enquote{} → hrvatski navodnici „..."
\usepackage{booktabs}
\usepackage{array}
\usepackage{listings}
\usepackage{xcolor}
\usepackage{colortbl}
\usepackage{hyperref}
\usepackage{tikz}
\usetikzlibrary{positioning,arrows.meta}

% ---------- Tipografija: IBM Plex ----------
\setsansfont{IBMPlexSans}[
  Extension      = .otf,
  UprightFont    = *-Regular,
  ItalicFont     = *-Italic,
  BoldFont       = *-SemiBold,
  BoldItalicFont = *-SemiBoldItalic
]
\setmainfont{IBMPlexSans}[
  Extension      = .otf,
  UprightFont    = *-Regular,
  ItalicFont     = *-Italic,
  BoldFont       = *-SemiBold,
  BoldItalicFont = *-SemiBoldItalic
]
\setmonofont{IBMPlexMono}[
  Extension      = .otf,
  UprightFont    = *-Regular,
  ItalicFont     = *-Italic,
  BoldFont       = *-SemiBold,
  BoldItalicFont = *-SemiBoldItalic,
  Scale          = 0.92
]

% ---------- Paleta ----------
\definecolor{ink}{HTML}{0C4685}      % primarna tamnoplava
\definecolor{ink2}{HTML}{6E7BAA}     % sekundarna plavo-siva
\definecolor{tint}{HTML}{EAEEF6}     % svjetli plavi tint
\definecolor{textclr}{HTML}{1A222E}  % glavni tekst
\definecolor{mutedclr}{HTML}{6B7484} % podnožje, sekundarni tekst
\definecolor{bgclr}{HTML}{F6F7F9}    % pozadina slajda
\definecolor{lineclr}{HTML}{D9DEE7}  % tanke linije tablica
\definecolor{divrule}{HTML}{9DB4D6}  % akcent na tamnim sekcijskim slajdovima
\definecolor{divtext}{HTML}{E6ECF6}  % tekst na tamnim sekcijskim slajdovima

\setbeamercolor{normal text}{fg=textclr,bg=bgclr}
\setbeamercolor{structure}{fg=ink}
\setbeamercolor{alerted text}{fg=ink}

\setbeamercolor{title}{fg=ink,bg=tint}
\setbeamercolor{subtitle}{fg=mutedclr}
\setbeamercolor{author}{fg=textclr}
\setbeamercolor{date}{fg=mutedclr}
\setbeamercolor{frametitle}{fg=ink,bg=tint}

\setbeamercolor{section in head/foot}{fg=white,bg=ink}
\setbeamercolor{subsection in head/foot}{fg=ink,bg=tint}
\setbeamercolor{author in head/foot}{fg=white,bg=ink}
\setbeamercolor{title in head/foot}{fg=textclr,bg=tint}
\setbeamercolor{date in head/foot}{fg=mutedclr,bg=tint}

\setbeamercolor{block title}{fg=white,bg=ink}
\setbeamercolor{block body}{fg=textclr,bg=tint}
\setbeamercolor{block title example}{fg=white,bg=ink2}
\setbeamercolor{block body example}{fg=textclr,bg=tint}
\setbeamercolor{itemize item}{fg=ink2}
\setbeamercolor{itemize subitem}{fg=ink2}
\setbeamercolor{enumerate item}{fg=ink2}

\setbeamertemplate{itemize items}[square]
\setbeamertemplate{navigation symbols}{}
\arrayrulecolor{lineclr}

% ---------- Kolegij, autori i logo ----------
\newcommand{\kolegij}{Programiranje za web}
\author{doc.\,dr.\,sc. Tomislav Petković \and mag.\,ing. Josip Torić}
\newcommand{\znak}[2]{\raisebox{-0.5\height}{\includegraphics[height=#1]{#2}}}
\titlegraphic{\znak{2.1cm}{../common/SIT_logo_tekst1.pdf}\hspace{1.8em}\znak{1.8cm}{../common/UNIZD_logo.pdf}}

% ---------- Zaglavlje ----------
% Jedna puna traka: lijevo sekcija, desno podsekcija.
\setbeamertemplate{headline}{%
  \leavevmode%
  \hbox{%
    \begin{beamercolorbox}[wd=\paperwidth,ht=2.25ex,dp=1ex,leftskip=.3cm,rightskip=.3cm]{section in head/foot}%
      \usebeamerfont{section in head/foot}\insertsectionhead\hfill\insertsubsectionhead%
    \end{beamercolorbox}%
  }%
}

% ---------- Podnožje ----------
% Lijevo: naziv trenutne prezentacije (kratki oblik naslova iz \title[...]{...})
% Sredina: naziv kolegija.  Desno: datum i redni broj slajda.
\setbeamertemplate{footline}{%
  \leavevmode%
  \hbox{%
    \begin{beamercolorbox}[wd=.34\paperwidth,ht=2.25ex,dp=1ex,center]{author in head/foot}%
      \usebeamerfont{author in head/foot}\insertshorttitle%
    \end{beamercolorbox}%
    \begin{beamercolorbox}[wd=.33\paperwidth,ht=2.25ex,dp=1ex,center]{title in head/foot}%
      \usebeamerfont{title in head/foot}\kolegij%
    \end{beamercolorbox}%
    \begin{beamercolorbox}[wd=.33\paperwidth,ht=2.25ex,dp=1ex,center]{date in head/foot}%
      \usebeamerfont{date in head/foot}\insertshortdate{}\hspace*{2em}%
      \insertframenumber{}\,/\,\inserttotalframenumber\hspace*{1ex}%
    \end{beamercolorbox}%
  }%
  \vskip0pt%
}

% ---------- Sekcijski slajd ----------
% \sekcija{Naslov sekcije}{jedna linija koja nabraja teme sekcije}
\newcommand{\sekcija}[2]{%
  \section{#1}%
  {%
    \setbeamercolor{background canvas}{bg=ink}%
    \setbeamertemplate{headline}{}%
    \begin{frame}
      \vfill
      \hspace*{0.03\textwidth}{\color{divrule}\rule{2.6cm}{1.1pt}}\par
      \vskip 0.7em
      \hspace*{0.03\textwidth}{\Large\bfseries\color{white}#1}\par
      \vskip 0.6em
      \hspace*{0.03\textwidth}\begin{minipage}{0.9\textwidth}\color{divtext}\small #2\end{minipage}\par
      \vfill
    \end{frame}%
  }%
}

% ---------- Listings ----------
\lstdefinelanguage{bash}{
  morekeywords={sudo,apt,brew,pip,pipx,uv,winget,python,git,flask,source,deactivate,code,gunicorn,docker,compose,heroku,render,railway,node,npm,pnpm,yarn,uvicorn,alembic,pytest,ruff,black,precommit},
  sensitive=false,
  morecomment=[l]{\#}
}
\lstdefinelanguage{python}{
  morekeywords={from,import,def,class,return,if,elif,else,for,in,while,try,except,raise,with,as,not,and,or,is,lambda,yield,True,False,None,await,async},
  sensitive=true,
  morecomment=[l]{\#},
  morestring=[b]',
  morestring=[b]"
}
\lstdefinelanguage{sql}{
  morekeywords={SELECT,FROM,WHERE,INSERT,INTO,VALUES,UPDATE,SET,DELETE,CREATE,TABLE,PRIMARY,KEY,FOREIGN,REFERENCES,UNIQUE,NOT,NULL,INDEX,ON,ALTER,DROP,CONSTRAINT},
  sensitive=false,
  morecomment=[l]{--}
}
\lstdefinelanguage{json}{
  morecomment=[l]{//},
  morestring=[b]",
  showstringspaces=false
}
\lstdefinelanguage{typescript}{
  morekeywords={import,export,from,const,let,var,function,return,if,else,for,of,in,while,switch,case,break,async,await,try,catch,finally,throw,new,class,this,extends,implements,interface,type,enum,readonly,public,private,protected,static,get,set,true,false,null,undefined,void,string,number,boolean,any,unknown,never,typeof,keyof,as},
  sensitive=true,
  morecomment=[l]{//},
  morecomment=[s]{/*}{*/},
  morestring=[b]',
  morestring=[b]",
  morestring=[b]`
}
\lstalias{javascript}{typescript}
\lstset{
  backgroundcolor=\color{tint},
  basicstyle=\ttfamily\footnotesize,
  keywordstyle=\color{ink}\bfseries,
  commentstyle=\color{mutedclr},
  frame=single,
  rulecolor=\color{lineclr},
  frameround=tttt,
  breaklines=true,
  tabsize=2,
  showstringspaces=false,
  upquote=true,
  columns=fullflexible,
  escapeinside={<@}{@>}
}
% ---------- Dijakritici u isječcima koda ----------
% listings čita bajt po bajt, pa višebajtni znak premjesti: "početni" ispadne
% "čpoetni". literate to ne popravlja. Rješenje je izlaz u LaTeX: sve između
% <@ i @> ne prolazi kroz listings, nego se slaže kao običan tekst, a unutar
% okvira je font ionako ttfamily, pa se ništa ne razlikuje.
%
%   async def get(...):
%       club = await repo.get_by_id(db, club_id)   <@\kom{dohvat kluba}@>
%       raise AppError("not_found", <@"Klub nije pronađen"@>, 404)
%
% Graničnici <@ i @> odabrani su jer se ne pojavljuju ni u jednom jeziku
% koji se na predavanjima citira; (* bi se sudaralo s def f(*args).
\newcommand{\kom}[1]{{\color{mutedclr}\# #1}}
% Isto za TypeScript, Vue i CSS, kojima komentar počinje s //:
%   const lifters = useLiftersStore()   <@\komts{dijeljeni popis}@>
\newcommand{\komts}[1]{{\color{mutedclr}// #1}}

% ---------- Činjenica naspram konvencije ----------
% Sve što nije propisano standardom ili jezikom, nego je dogovor struke,
% ide u ovaj okvir, da student zna koje je vrste tvrdnja pred njim.
% \begin{konvencija} ... \end{konvencija}
% \begin{konvencija}[Konvencija: nazivi putanja] ... \end{konvencija}
\newenvironment{konvencija}[1][Konvencija]
  {\begin{exampleblock}{#1}}
  {\end{exampleblock}}

% ---------- Helper macros ----------
% Stupac tablice zadane širine, poravnat ulijevo (bez rastezanja i dijeljenja riječi).
\newcolumntype{L}[1]{>{\raggedright\arraybackslash}p{#1}}
\newcommand{\codefile}[1]{\texttt{#1}}
\newcommand{\kbd}[1]{\texttt{#1}}

% ---------- Code frame environment ----------
% Automatski dodaje fragile i shrink=5, koristi se umjesto \begin{frame}
% Primjer: \begin{codeframe}{Naslov} ... \end{codeframe}
% Primjer s opcijama: \begin{codeframe}[shrink=15]{Naslov} ... \end{codeframe}
\newenvironment{codeframe}[2][]
  {\begin{frame}[fragile,shrink=5,#1]{#2}}
  {\end{frame}}


\title[Predavanje 9: Upravljanje stanjem]{Programiranje za web \\ Predavanje 9: Upravljanje stanjem i stanja sučelja}
\date{10.12.2026.}

\begin{document}

\frame{\titlepage}

% ------------------------------------------------------------
\begin{frame}{Plan predavanja}
\begin{enumerate}
  \item Pohrana koju dijele prikazi
  \item Stanja jednog prikaza
  \item Obrasci i odbijeni zahtjevi
  \item Ono što se ponavlja
\end{enumerate}

\vspace{0.6em}
Zahtjevi odlaze i odgovori stižu. Ovo predavanje uređuje gdje odgovor stoji dok je aplikacija otvorena i što korisnik vidi dok ga nema.
\end{frame}

% ------------------------------------------------------------
\begin{frame}{Ishodi učenja}
Nakon ovog predavanja student može:

\begin{itemize}
  \item odlučiti pripada li podatak jednoj komponenti ili pohrani koju dijele prikazi
  \item napisati pohranu sa stanjem, izvedenim vrijednostima i radnjama
  \item prikazati učitavanje, pogrešku i prazan skup rezultata
  \item povezati obrazac s podatkom i poslati ga
  \item prikazati odgovor \texttt{422} uz polje, a \texttt{400} i \texttt{409} iznad obrasca
  \item izdvojiti dio prikaza ili logike koji se ponavlja
\end{itemize}
\end{frame}

% ------------------------------------------------------------
\begin{frame}{Domena primjera}
Klub u aplikaciji vodi svoje natjecatelje i njihove prijave. Isti popis natjecatelja treba na nekoliko mjesta: u pregledu kluba, u obrascu za prijavu na natjecanje i u navigaciji, kao broj uz stavku.

\vspace{0.9em}
Backend svaki od tih prikaza poslužuje istom rutom. Pitanje je koliko puta aplikacija tu rutu poziva i gdje odgovor čuva dok je otvorena.

\vspace{0.9em}
Uz to, svaki obrazac koji klub pošalje poslužitelj može odbiti na tri načina, a svaki od njih traži drukčiji prikaz.
\end{frame}

% ------------------------------------------------------------
\begin{frame}[fragile]{Podatak koji treba dvama prikazima (I)}
Natjecatelje dohvaća prikaz koji ih nabraja:

\begin{lstlisting}[language=typescript]
const lifters = ref<Lifter[]>([])
onMounted(async () => {
  lifters.value = await getLifters(clubId)
})
\end{lstlisting}

\vspace{0.5em}
Obrascu za prijavu na natjecanje treba isti popis, da ponudi izbor natjecatelja. Napisan na isti način, on isti zahtjev šalje još jednom.

\vspace{0.5em}
Natjecatelj dodan u jednom prikazu u drugome se ne pojavljuje dok se taj prikaz ponovno ne otvori.
\end{frame}

% ------------------------------------------------------------
\begin{frame}{Podatak koji treba dvama prikazima (II)}
Prvi odgovor je proslijediti popis kroz svojstva komponenata. On vrijedi dok su prikazi u istom stablu i dok ih je malo.

\vspace{0.9em}
Popis koji treba navigaciji, pregledu kluba i obrascu nije u istom stablu ni na istoj razini. Prosljeđivanje bi tada išlo kroz komponente koje taj popis ne prikazuju, samo da bi ga predale dalje.

\vspace{0.9em}
Drugi je odgovor podatak izdvojiti iz stabla komponenata: držati ga na jednom mjestu izvan njih, odakle ga svaki prikaz čita izravno.
\end{frame}

% ============================================================
\sekcija{Pohrana koju dijele prikazi}{Što pohrana drži, kako se piše, zamka pri rastavljanju i granica prema lokalnim vrijednostima.}

% ------------------------------------------------------------
\begin{frame}{Što pohrana drži}
Pohrana (engl. \emph{store}) je objekt koji postoji dok je aplikacija otvorena i ne pripada nijednoj komponenti. Ima tri vrste sadržaja:

\vspace{0.6em}
\begin{center}
{\small
\setlength{\tabcolsep}{6pt}
\renewcommand{\arraystretch}{1.35}
\begin{tabular}{@{}L{2.9cm}L{3.5cm}L{4.7cm}@{}}
  \toprule
  \textbf{Sadržaj} & \textbf{Zapisano kao} & \textbf{Odgovara sloju} \\
  \midrule
  stanje & \texttt{ref()} & podatci u bazi \\
  izvedene vrijednosti & \texttt{computed()} & izračun nad podatcima \\
  radnje & obične funkcije & \codefile{services} \\
  \bottomrule
\end{tabular}
}
\end{center}

\vspace{0.6em}
Podjela je ista kao na poslužitelju: ono što jest, ono što se iz toga računa i ono što to mijenja. Razlika je u tome što pohrana nestaje pri osvježavanju stranice.
\end{frame}

% ------------------------------------------------------------
\begin{frame}[fragile]{Pohrana za natjecatelje}
\vspace{-0.4em}
\begin{lstlisting}[language=typescript]
// stores/lifters.ts
export const useLifterStore = defineStore('lifters', () => {
  const lifters = ref<Lifter[]>([])
  const ucitavanje = ref(false)
  const greska = ref<string | null>(null)
  async function dohvati(clubId: number) {
    ucitavanje.value = true
    greska.value = null
    try {
      lifters.value = await getLifters(clubId)
    } catch {
      greska.value = 'Popis natjecatelja nije dohvacen'
    }
    ucitavanje.value = false
  }
  return { lifters, ucitavanje, greska, dohvati }
})
\end{lstlisting}
\end{frame}

% ------------------------------------------------------------
\begin{frame}[fragile]{Pohrana u komponenti}
\begin{lstlisting}[language=typescript]
const store = useLifterStore()
const { lifters } = store            <@\kom{prikaz se vise ne osvjezava}@>

import { storeToRefs } from 'pinia'
const { lifters, ucitavanje } = storeToRefs(store)
const { dohvati } = store            <@\kom{funkcije smiju izravno}@>
\end{lstlisting}

\vspace{0.5em}
Rastavljanje objekta kopira trenutnu vrijednost. Kopija nije praćena, pa se prikaz koji je čita prestane osvježavati kad se pohrana promijeni, i to bez ikakve poruke.

\vspace{0.5em}
\texttt{storeToRefs} umjesto vrijednosti vraća praćene reference. Funkcije se smiju uzeti izravno, jer one nisu praćene vrijednosti.
\end{frame}

% ------------------------------------------------------------
\begin{frame}{Pohrana ili lokalna vrijednost}
\vspace{0.3em}
\begin{center}
{\small
\setlength{\tabcolsep}{6pt}
\renewcommand{\arraystretch}{1.3}
\begin{tabular}{@{}L{2.9cm}L{4.0cm}L{4.2cm}@{}}
  \toprule
   & \textbf{Lokalna vrijednost} & \textbf{Pohrana} \\
  \midrule
  Dopire do & jedne komponente & svakog prikaza \\
  Traje & dok je komponenta na zaslonu & dok je aplikacija otvorena \\
  Primjer & upisani tekst, otvoren izbornik & prijavljeni korisnik, popis natjecatelja \\
  \bottomrule
\end{tabular}
}
\end{center}

\vspace{0.6em}
Podatak koji trebaju dva prikaza ide u pohranu. Podatak koji postoji samo dok je jedan obrazac otvoren ostaje u komponenti.

\vspace{0.4em}
Sve u pohranu nije rješenje. Pohrana koja drži i privremene vrijednosti jednog obrasca zadržava ih i nakon što korisnik ode s te stranice.
\end{frame}

% ------------------------------------------------------------
\begin{frame}{Vježba: isti podatak u dvama prikazima}
Prijavljeni ste kao klub \texttt{behemot@plreg.hr} i imate prikaz koji nabraja natjecatelje.

\vspace{0.8em}
\begin{enumerate}
  \itemsep0.4em
  \item Dodajte drugi prikaz koji ispisuje samo broj natjecatelja i dohvaća ih na isti način.
  \item Otvorite karticu Mreža i prijeđite s jednog prikaza na drugi i natrag.
  \item Prenesite popis u pohranu i neka oba prikaza čitaju nju, a dohvat pokreće samo prvi.
  \item Ponovite drugi korak.
  \item Rastavite pohranu bez \texttt{storeToRefs} i dodajte natjecatelja.
\end{enumerate}

\vspace{0.8em}
Koliko je zahtjeva poslano prije, a koliko poslije trećeg koraka, i što se dogodilo u petom?
\end{frame}

% ============================================================
\sekcija{Stanja jednog prikaza}{Prikaz dok odgovora nema, prazan skup rezultata, pogreška i kada ponovni pokušaj ima smisla.}

% ------------------------------------------------------------
\begin{frame}[fragile]{Prikaz koji pokazuje samo podatke}
Prikaz napisan tako da ispisuje popis i ništa drugo:

\begin{lstlisting}[language=html]
<table>
  <tr v-for="l in store.lifters" :key="l.id">
    <td>{{ l.last_name }}, {{ l.first_name }}</td>
  </tr>
</table>
\end{lstlisting}

\vspace{0.5em}
Dok odgovor putuje, tablica je prazna. Ako zahtjev ne prođe, tablica je prazna. Ako klub nema nijednog natjecatelja, tablica je prazna.

\vspace{0.5em}
Tri različita stanja izgledaju isto, pa korisnik ne zna čeka li, je li zahtjev odbijen ili doista nema podataka.
\end{frame}

% ------------------------------------------------------------
\begin{frame}{Tri stanja svakog prikaza}
\vspace{0.3em}
\begin{center}
{\small
\setlength{\tabcolsep}{6pt}
\renewcommand{\arraystretch}{1.35}
\begin{tabular}{@{}L{2.6cm}L{3.4cm}L{5.1cm}@{}}
  \toprule
  \textbf{Stanje} & \textbf{Uvjet} & \textbf{Što se prikazuje} \\
  \midrule
  učitavanje & odgovor još nije stigao & znak da se čeka \\
  pogreška & zahtjev nije prošao & poruka i, kad ima smisla, ponovni pokušaj \\
  prazno & odgovor je prazan popis & rečenica i put do prvog upisa \\
  \bottomrule
\end{tabular}
}
\end{center}

\vspace{0.6em}
Provjere se nižu tim redom: prvo učitavanje, pa pogreška, pa prazno, pa podatci. Obrnut redoslijed pokaže prazan popis dok odgovor još putuje.

\vspace{0.4em}
Prazan skup rezultata svaki novi korisnik vidi prvi, pa je to mjesto na kojem mu se kaže što dalje.
\end{frame}

% ------------------------------------------------------------
\begin{frame}[fragile]{Ista četiri slučaja u prikazu}
\begin{lstlisting}[language=html]
<div v-if="ucitavanje">Ucitavanje...</div>

<div v-else-if="greska" class="error">
  {{ greska }}
  <button @click="dohvati(clubId)">Pokusaj ponovno</button>
</div>

<p v-else-if="lifters.length === 0">
  Nemate upisanih natjecatelja.
  <RouterLink to="/lifters/new">Upisite prvoga.</RouterLink>
</p>

<table v-else> ... </table>
\end{lstlisting}

\vspace{0.4em}
Niz \texttt{v-if} i \texttt{v-else-if} jamči da je u svakom trenutku prikazano točno jedno od toga.
\end{frame}

% ------------------------------------------------------------
\begin{frame}{Kada ponovni pokušaj ima smisla}
\vspace{0.3em}
\begin{center}
{\small
\setlength{\tabcolsep}{6pt}
\renewcommand{\arraystretch}{1.3}
\begin{tabular}{@{}L{3.2cm}L{4.0cm}L{3.9cm}@{}}
  \toprule
  \textbf{Uzrok} & \textbf{Primjer} & \textbf{Ponovni pokušaj} \\
  \midrule
  mreža & veza je prekinuta & ima smisla \\
  poslužitelj, \texttt{5xx} & pogreška u obradi & ima smisla \\
  zahtjev, \texttt{4xx} & podatci ne prolaze provjeru & nema smisla \\
  \bottomrule
\end{tabular}
}
\end{center}

\vspace{0.6em}
Isti zahtjev koji je odbijen zbog svojeg sadržaja bit će odbijen i drugi put. Gumb za ponovni pokušaj uz takvu poruku obećava nešto što se neće dogoditi.

\vspace{0.4em}
Odgovore \texttt{401} i \texttt{403} prikaz uopće ne vidi: njih obrađuje presretač.
\end{frame}

% ------------------------------------------------------------
\begin{frame}{Vježba: tri stanja}
Radite nad prikazom koji nabraja natjecatelje.

\vspace{0.8em}
\begin{enumerate}
  \itemsep0.4em
  \item Dodajte sva tri stanja i poredajte provjere kako je opisano.
  \item Zaustavite backend i otvorite prikaz.
  \item Pokrenite backend, prijavite se kao klub \texttt{sokol@plreg.hr} i otvorite isti prikaz.
  \item U alatima za razvoj usporite mrežu i ponovno otvorite prikaz.
  \item Zamijenite redoslijed tako da se prazan popis provjerava prvi i ponovite četvrti korak.
\end{enumerate}

\vspace{0.8em}
Koje ste stanje vidjeli u svakom od koraka 2, 3 i 4, a što se promijenilo u petom?
\end{frame}

% ============================================================
\sekcija{Obrasci i odbijeni zahtjevi}{Obrazac vezan uz podatak, provjera prije slanja, tri oblika odbijanja i mjesto na kojem se poruka prikazuje.}

% ------------------------------------------------------------
\begin{frame}[fragile]{Obrazac i podatak}
Obrazac se ne čita iz stabla elemenata. Polja su vezana uz jedan objekt, a on je ono što se šalje:

\begin{lstlisting}[language=html]
<script setup lang="ts">
const obrazac = ref({
  first_name: '', last_name: '', birth_date: '', gender: 'M',
})
</script>

<template>
  <input v-model="obrazac.first_name" required>
  <input v-model="obrazac.birth_date" type="date" required>
</template>
\end{lstlisting}

\vspace{0.4em}
Polja odgovaraju shemi unosa na poslužitelju, jedno po jedno. Oznaku kluba obrazac ne nosi, jer ona stoji u adresi rute.
\end{frame}

% ------------------------------------------------------------
\begin{frame}{Provjera prije slanja}
Preglednik sam odbija prazno polje označeno s \texttt{required} i vrijednost koja ne odgovara tipu polja. To se dogodi bez zahtjeva, pa je odgovor trenutačan.

\vspace{0.9em}
Ta provjera postoji radi korisnika. Zahtjev poslan mimo obrasca je ne vidi, pa ista pravila stoje i u shemi na poslužitelju.

\vspace{0.9em}
Zato provjera u pregledniku ne mora biti potpuna. Ima smisla ondje gdje štedi zahtjev i gdje je poruka brža: prazna polja, očit oblik datuma, najmanja i najveća duljina.
\end{frame}

% ------------------------------------------------------------
\begin{frame}[fragile]{Tri oblika odbijanja}
Isti obrazac poslužitelj odbija na tri načina, a oblik odgovora je uvijek isti:

\begin{lstlisting}[language=json]
{"code": "validation_error", "message": "Neispravni podatci u zahtjevu",
 "errors": [{"field": "birth_date",
             "message": "Datum rodenja mora biti u proslosti"}]}
\end{lstlisting}

\vspace{0.4em}
\begin{center}
{\small
\setlength{\tabcolsep}{6pt}
\renewcommand{\arraystretch}{1.3}
\begin{tabular}{@{}L{1.4cm}L{4.6cm}L{5.1cm}@{}}
  \toprule
  \textbf{Kod} & \textbf{Uzrok} & \textbf{Gdje se prikazuje} \\
  \midrule
  \texttt{422} & polje ne odgovara shemi & uz to polje \\
  \texttt{409} & zapis bi prekršio jedinstvenost & iznad obrasca \\
  \texttt{400} & poslovno pravilo ne dopušta & iznad obrasca \\
  \bottomrule
\end{tabular}
}
\end{center}
\end{frame}

% ------------------------------------------------------------
\begin{frame}[fragile]{Poruka uz polje}
Samo odgovor \texttt{422} nosi polje \texttt{errors}, a u njemu ime polja na koje se poruka odnosi:

\begin{lstlisting}[language=typescript]
// utils/greske.ts
export function poPolju(data: ApiError): Record<string, string> {
  const izlaz: Record<string, string> = {}
  for (const g of data.errors ?? []) izlaz[g.field] = g.message
  return izlaz
}
\end{lstlisting}

\vspace{0.4em}
Rezultat je objekt koji za svako ime polja daje poruku, pa ga predložak čita izravno:

\begin{lstlisting}[language=html]
<span v-if="greske.birth_date" class="error">
  {{ greske.birth_date }}
</span>
\end{lstlisting}
\end{frame}

% ------------------------------------------------------------
\begin{frame}[fragile]{Slanje obrasca}
\begin{lstlisting}[language=typescript]
async function posalji() {
  greske.value = {}
  poruka.value = ''
  try {
    await api.post(`/clubs/${clubId}/lifters`, obrazac.value)
    router.push('/lifters')
  } catch (e) {
    const odgovor = (e as AxiosError<ApiError>).response
    if (odgovor?.status === 422) greske.value = poPolju(odgovor.data)
    else poruka.value = odgovor?.data.message ?? 'Zahtjev nije prosao'
  }
}
\end{lstlisting}

\vspace{0.4em}
Polje \texttt{message} namijenjeno je čovjeku, pa se prikazuje onakvo kakvo je. Grananje ide po \texttt{code}, ne po tekstu poruke.
\end{frame}

% ------------------------------------------------------------
\begin{frame}[fragile]{Obrazac za izmjenu}
Obrazac za izmjenu počinje popunjen, pa prvo dohvaća zapis:

\begin{lstlisting}[language=typescript]
async function ucitaj(id: number) {
  const { data } = await api.get<Lifter>(
    `/clubs/${clubId}/lifters/${id}`)
  obrazac.value = { ...data }
}

onMounted(() => ucitaj(Number(route.params.lifterId)))
watch(() => route.params.lifterId, (id) => ucitaj(Number(id)))
\end{lstlisting}

\vspace{0.4em}
\texttt{onMounted} pokriva prvi dolazak na rutu. Promatrač pokriva prijelaz s jednog natjecatelja na drugoga: adresa se promijeni, a komponenta ostaje ista, pa bi bez njega obrazac zadržao prethodni zapis.
\end{frame}

% ------------------------------------------------------------
\begin{frame}{Vježba: tri odbijena obrasca}
Prijavljeni ste kao klub \texttt{behemot@plreg.hr}, a \emph{Zimski kup} je u prvom razdoblju.

\vspace{0.8em}
\begin{enumerate}
  \itemsep0.4em
  \item U obrascu za novog natjecatelja upišite datum rođenja u budućnosti i pošaljite.
  \item Prijavite natjecatelja na natjecanje, pa istoga natjecatelja prijavite još jednom.
  \item U obrascu za prijavu upišite kategoriju koja ne odgovara spolu natjecatelja.
  \item Kao organizator \texttt{admin@plreg.hr} pomaknite preliminarni rok u prošlost, pa se vratite kao klub i ponovite drugi korak s drugim natjecateljem.
\end{enumerate}

\vspace{0.6em}
Koji ste \texttt{code} dobili u svakom koraku i gdje bi u obrascu svaka od tih poruka trebala stajati?
\end{frame}

% ============================================================
\sekcija{Ono što se ponavlja}{Isti dio prikaza na više mjesta, komponenta s vlastitim sučeljem, ista logika u više komponenata.}

% ------------------------------------------------------------
\begin{frame}{Isti dio prikaza na više mjesta}
Obrazac za natjecatelja ima četiri polja. Svako od njih je oznaka, polje za unos i mjesto za poruku o pogrešci, dakle isti raspored četiri puta.

\vspace{0.9em}
Obrazac za prijavu na natjecanje ima iste takve dijelove. Isto vrijedi za svaki sljedeći obrazac.

\vspace{0.9em}
Kad se raspored prepiše, mijenja se na svakom mjestu posebno. Ako se na jednom od njih poruka o pogrešci zaboravi, obrazac radi, a korisnik ne vidi zašto zahtjev ne prolazi.
\end{frame}

% ------------------------------------------------------------
\begin{frame}[fragile]{Polje obrasca kao komponenta}
\begin{lstlisting}[language=html]
<!-- components/PoljeObrasca.vue -->
<script setup lang="ts">
defineProps<{ oznaka: string; greska?: string;
              modelValue: string }>()
defineEmits<{ 'update:modelValue': [vrijednost: string] }>()
</script>
<template>
  <label>{{ oznaka }}</label>
  <input :value="modelValue"
    @input="$emit('update:modelValue',
                  ($event.target as HTMLInputElement).value)">
  <span v-if="greska" class="error">{{ greska }}</span>
</template>
\end{lstlisting}

\vspace{0.4em}
Par \texttt{modelValue} i \texttt{update:modelValue} dogovoreno je ime koje \texttt{v-model} prepoznaje, pa se komponenta koristi kao obično polje: \texttt{<PoljeObrasca v-model="obrazac.first\_name" />}.
\end{frame}

% ------------------------------------------------------------
\begin{frame}[fragile]{Oznaka stanja prijave}
Prijava ima stanje \texttt{active} ili \texttt{withdrawn}, a popis ga prikazuje na svakom retku:

\begin{lstlisting}[language=html]
<!-- components/OznakaStanja.vue -->
<script setup lang="ts">
const props = defineProps<{ stanje: string }>()
const razred = computed(() =>
  props.stanje === 'active' ? 'badge-green' : 'badge-grey')
</script>

<template>
  <span :class="razred">{{ stanje }}</span>
</template>
\end{lstlisting}

\vspace{0.4em}
Komponenta ne odlučuje što se s prijavom smije: gumb za odjavu stoji u popisu, uz uvjet koji ovisi o stanju prijave i razdoblju natjecanja.
\end{frame}

% ------------------------------------------------------------
\begin{frame}{Kada izdvojiti, a kada ne}
Izdvaja se ono što se ponavlja i ima jasnu granicu: polje obrasca, poruka o pogrešci s ponovnim pokušajem, oznaka stanja, zaglavlje tablice.

\vspace{0.9em}
Ne izdvaja se ono što postoji na jednom mjestu. Komponenta koja se koristi jednom dodaje datoteku, a ne uklanja ponavljanje.

\vspace{0.9em}
Komponenta s desetak svojstava znak je da radi previše. Takva se rastavlja ili joj se dio sadržaja predaje kao umetak, čime roditelj odlučuje što se prikazuje, a komponenta gdje.
\end{frame}

% ------------------------------------------------------------
\begin{frame}[fragile]{Ista logika u više komponenata}
Ponavlja se i ono što nije prikaz. Odgoda slanja pri tipkanju treba svakom polju za pretragu:

\begin{lstlisting}[language=typescript]
// composables/useOdgoda.ts
export function useOdgoda(kasnjenje = 300) {
  const unos = ref('')
  const odgodeni = ref('')
  let mjerac: ReturnType<typeof setTimeout>
  watch(unos, (v) => {
    clearTimeout(mjerac)
    mjerac = setTimeout(() => { odgodeni.value = v }, kasnjenje)
  })
  return { unos, odgodeni }
}
\end{lstlisting}

\vspace{0.4em}
Bez odgode pretraga po prezimenu šalje zahtjev na svaki pritisnut znak. Uz nju se šalje jedan, kad korisnik prestane tipkati.
\end{frame}

% ------------------------------------------------------------
\begin{frame}{Funkcija koja vraća praćene vrijednosti}
Takva se funkcija po dogovoru imenuje s predmetkom \texttt{use} i vraća praćene vrijednosti umjesto običnih.

\vspace{0.9em}
Razlika prema pohrani je u tome tko je vlasnik. Pohrana postoji jednom za cijelu aplikaciju, a ovakva funkcija svakom pozivatelju daje vlastiti skup vrijednosti.

\vspace{0.9em}
Zato prijavljeni korisnik stoji u pohrani, a odgoda unosa u ovakvoj funkciji: dva polja za pretragu na istoj stranici moraju imati odvojene mjerače.
\end{frame}

% ------------------------------------------------------------
\begin{frame}{Vježba: izdvojena komponenta}
Imate obrazac za natjecatelja s četiri polja i prikazom poruka iz odgovora \texttt{422}.

\vspace{0.8em}
\begin{enumerate}
  \itemsep0.4em
  \item Izdvojite polje obrasca u komponentu i upotrijebite je za sva četiri polja.
  \item Pošaljite obrazac s praznim prezimenom i datumom rođenja u budućnosti.
  \item Dodajte komponenti svojstvo kojim se polje označava kao obavezno i upotrijebite ga.
  \item Upotrijebite istu komponentu u obrascu za prijavu na natjecanje.
\end{enumerate}

\vspace{0.8em}
Koliko ste mjesta morali promijeniti u trećem koraku, a koliko biste ih promijenili da polja nisu izdvojena?
\end{frame}

% ------------------------------------------------------------
\begin{frame}{Sažetak}
\begin{itemize}
  \itemsep0.3em
  \item Podatak koji trebaju dva prikaza stoji izvan stabla komponenata.
  \item Pohrana ima stanje, izvedene vrijednosti i radnje, kao i sloj na poslužitelju.
  \item Rastavljanje pohrane bez \texttt{storeToRefs} tiho prekida praćenje.
  \item Privremene vrijednosti jednog obrasca ostaju u komponenti.
  \item Prikaz bez triju stanja izgleda jednako u sva tri.
  \item Redoslijed provjere je učitavanje, pogreška, prazno, podatci.
  \item Ponovni pokušaj ima smisla samo za mrežu i \texttt{5xx}.
  \item Provjera u pregledniku je radi brzine, shema na poslužitelju radi točnosti.
  \item \texttt{422} se prikazuje uz polje, \texttt{409} i \texttt{400} iznad obrasca.
  \item Pohrana je zajednička, a \texttt{use} funkcija daje svakom pozivatelju svoju.
\end{itemize}
\end{frame}

\end{document}
